\subsection{How Do We Build a Virtuous Cycle?}

The short answer is to make the next good action easier than the next bad
one, then repeat that advantage until the new pattern starts to carry itself.
A virtuous cycle is not built by one large burst of willpower. It is built by
arranging the environment, timing, and first step so that each successful
transition lowers the next barrier and slightly strengthens the target state.
In the language of this book, the goal is to reduce framework inertia at the
moment of entry, spend volitional energy where it has the highest leverage,
and keep the decision window short enough that the old frame cannot fully
reassert itself.

\begin{enumerate}[nosep]
  \item Name one concrete target state in plain language. Do not start with a
        vague identity claim; start with the behavior you want to make easier.
  \item Identify the smallest decision window in which the transition can
        still happen. If the window is long, shrink it by pre-committing the
        first action.
  \item Remove one barrier variable before the next attempt. Move objects,
        prepare the workspace, reduce ambiguity, or eliminate a predictable
        interruption.
  \item Define one minimum viable entry action that is so small it can happen
        even when volitional energy is only moderate.
  \item Repeat the same entry condition several times so the target state
        begins to feel familiar and less expensive to enter.
  \item After each success, make one local improvement that reduces the next
        start cost, not the next abstract aspiration.
\end{enumerate}

\begin{remark}
A student who wants a virtuous study cycle should not begin by promising to
``be disciplined.'' A better move is to place the book open on the desk, write
the first problem number in advance, and decide that the session starts with
five minutes of work. The first success lowers the next barrier, and the
routine becomes easier because the system now expects the target state.
\end{remark}
